\section{模型范围与符号说明}
\label{sec:symbols}

\subsection{系统边界与时间尺度}

模型的空间边界从风、光、潮等发电单元的设备端开始，经海上集电系统和公共交流母线，
延伸至四类结算计量点：海缆陆上受端、氢气管输或船运接收点、算力服务结算计量点以及海洋用能
设备入口。源侧集电损耗、海缆损耗、电解系统辅助用电和数据中心辅助用电分别在所属模块
中计算，公共母线只汇总各模块的最终功率，避免重复扣减。

调度时段集合及主要设备集合定义为
\begin{equation}
\mathcal T=\{1,2,\ldots,T\},\qquad
\mathcal S=\{w,pv,tc\},\qquad
\mathcal J=\{E,H,C,M\},\qquad
\Delta t=1\myunit{h}.
\label{eq:common-sets}
\end{equation}
其中，$\mathcal S$分别表示风电、光伏和潮流能，$\mathcal J$分别表示电力、氢能、算力和
海洋用能路径。功率变量表示时段平均值，状态变量取时段末值并传递至下一时段。

\subsection{模块划分与接口}

各子模型的输入输出见表~\ref{tab:module-responsibility}。联合调度模型只读取这些接口量，
不在公共层重复建立设备内部方程。

\begin{table}[H]
\centering
\caption{第四章各子模型的主要输入与输出}
\label{tab:module-responsibility}
\begin{tabularx}{\textwidth}{>{\centering\arraybackslash}p{1.2cm} p{4.0cm} X}
\toprule
小节 & 模型内容 & 向联合调度提供的变量 \\
\midrule
4.2 & 多源可再生能源出力 & 逐源可用功率、实际利用功率、弃电量和设备状态 \\
4.3 & 产品输出通道 & 海缆受端电量、氢气接收量、有效算力服务量和海洋用能量 \\
4.4 & 电池储能 & 充放电功率、时段末能量和构网备用裕度 \\
4.5 & 制氢与储氢 & 电解功率、产氢量、氢库存、氢气供应量和可选氢转电功率 \\
4.6 & 绿色算力 & 任务完成量、任务积压、IT功率、设施功率、PUE和SLA指标 \\
4.7 & 公共母线平衡 & 各时段功率平衡残差和损耗核算结果 \\
4.8 & 目标函数 & 运行成本、收入、弃电、可靠性和环境评价指标 \\
4.9--4.10 & 联合约束与求解 & 调度动作、状态轨迹、求解状态和数值检查结果 \\
\bottomrule
\end{tabularx}
\end{table}

\subsection{单位与符号约定}

本报告将物理电量与算力服务量分开计量。功率统一采用MW，电能采用MWh，氢气质量采用kg。
算力任务采用标准化算力服务量（compute service，CS）表示：任务到达量、完成量和队列量
采用MWh-CS，服务速率采用MW-CS。MWh-CS是面向调度的参考等效工作量，并非电能单位，
也不直接等同于GPU$\cdot$h、Token或任务数。

IT设备的实际功率和电量分别采用MW-IT和MWh-IT；数据中心设施入口的总功率和总电量采用
MW和MWh。三者之间的关系在4.6中通过负载率、IT负载--功率模型和PUE模型建立。全文均按上述
定义使用符号：MWh-CS只描述任务服务量，MWh-IT只描述IT设备耗电。

\begin{table}[H]
\centering
\caption{第四章主要计量单位}
\label{tab:unit-convention}
\begin{tabularx}{\textwidth}{p{3.0cm}p{2.2cm}X}
\toprule
对象 & 单位 & 含义 \\
\midrule
设备功率、母线功率 & MW & 时段平均有功功率 \\
电池能量、设施耗电 & MWh & 物理电能 \\
算力服务速率 & MW-CS & 单位时间内完成的标准化算力工作量 \\
算力服务量、任务队列 & MWh-CS & 标准化算力工作量 \\
IT设备功率、IT耗电 & MW-IT、MWh-IT & IT设备侧实际功率和电能 \\
氢气产量与库存 & kg & 时段产氢量、供应量和库存 \\
\bottomrule
\end{tabularx}
\end{table}

\subsection{状态传递、数据来源与参数处理}

电池能量$E_t^B$、氢库存$H_t$、电解槽在线模块数、设备状态和柔性算力队列$Q_t$均按
时段递推。第$t$时段的决策使用期初状态以及当时可获得的资源、需求、价格和事件信息；
第$t+1$时段的实际出力与价格不进入当前决策。输入数据采用统一时间戳、单位和缺失值处理
规则，求解后分别检查功率平衡、能量递推和质量守恒残差。

模型参数按数据基础分为三类：第一类为标准、设备手册或公开试验资料能够支持的技术参数；
第二类为NASA POWER、Copernicus Marine等公开数据形成的环境输入；第三类为研究者为构造
基准算例所设定的容量、价格、故障和控制参数。第三类参数仅用于机制分析与方案比较，均在
第五章参数表中列明，并通过敏感性分析考察其影响。工程应用阶段需采用场址观测、OEM曲线、
设计计算和合同数据更新相应参数。
